\subsection{Rank-one modules and semisimple Hopf algebras}
\begin{proposition}\label{prop:hopf-small}
The rings $R8\_2\_41$, $R8\_2\_86$, $R8\_2\_162$, $R8\_2\_140$ and
$R8\_2\_784$ have no complex fusion categorification.
\end{proposition}
\begin{proof}
For the first three rings, a faithful $C_2$ grading has a rank-one
nonneutral component and a dimension-sixteen neutral category
$\mathcal D$ of type $(1^4,2^3)$. A rank-one module supplies a fiber
functor, so $\mathcal D=\Rep(H)$ for a semisimple Hopf algebra of
 dimension sixteen. Invertibility of the component further requires
$\Rep(H^*)\simeq\mathcal D$; in particular both groups of grouplikes
have order four. Kashina's classification
\cite[Theorem~1.2 and Table~1]{Kashina} leaves only $H_{C:\sigma_1}$,
whose representation ring is $\Rep(D_{16})$: its invertibles form
$C_2^2$, and all three dimension-two simples are self-dual.
The rings $R8\_2\_41$, $R8\_2\_86$ instead have cyclic invertibles,
while $R8\_2\_162$ has a nonselfdual dimension-two pair.

For $R8\_2\_140$, census indices $\{0,1,2,3,7\}$ give the neutral
Tambara--Yamagami category on $C_4$, of dimension eight. Each other
simple separately gives a rank-one left module: an invertible acts as
one and the dimension-two object acts as two. Thus the neutral category
has a fiber functor and is $\Rep(H)$ with $G(H^*)=C_4$.
Kashina's prime-power theorem \cite[Theorem~1.1]{Kashina} says that a
semisimple Hopf algebra of prime-power dimension with cyclic grouplikes
is a group algebra. It would give $H^*\simeq\C[C_8]$, which has eight
grouplikes, a contradiction. This argument does not require the
rank-one modules to be grading components.

For $R8\_2\_784$, the neutral component has dimension twenty-two,
rank seven and type $(1^2,2^5)$; the other component has a unique
simple of dimension $\sqrt{22}$. Its fiber functor reconstructs a
Hopf algebra of dimension $2\cdot11$. The Hopf $pq$ classification
\cite{HopfPQ} and the nonpointed neutral rules force
$\mathcal D=\Rep(D_{22})$. Its fiber-functor dual is
$\Vec_{D_{22}}$, of rank twenty-two. The invertible grading component
would identify this dual with $\mathcal D$, of rank seven. This is
impossible.
\end{proof}

\subsection{Preserved Tannakian subcategories}
\begin{proposition}\label{prop:prime-r8}
The rings $R8\_2\_697$, $R8\_2\_808$ and $R8\_2\_181$ have no
complex fusion categorification.
\end{proposition}
\begin{proof}
For $R8\_2\_697$, dimension sixty-eight, type $(1^4,4^4)$ and trivial
universal grading, with all four invertibles fixing each noninvertible,
force a $C_2$ equivariantization with descent
$\mathcal D$ of dimension thirty-four and type $(1^2,2^8)$.
Its adjoint is the whole category: a proper adjoint has dimension
$1,2$ or $17$; a prime-dimensional adjoint is pointed, but these
possibilities cannot contain the required dimension-four squares with
only two invertibles. Solvability and prime descent now give
$\mathcal D=(\Vec_{C_{17}}^\omega)^{C_2}$, with inversion.
Since $U(\mathcal D)=1$, central lifts of its invertibles are unique.
The outer involution therefore preserves its central $\Rep(C_2)$,
and the combined equivariantization group is $C_2^2$, as its four
linear characters are the census invertibles. Every
$C_2^2\to\operatorname{Aut}(C_{17})$ has image of order at most two.
A nonzero orbit has size two and cyclic stabilizer; it therefore
supplies dimension-two equivariant simples. The target has only
dimension-four noninvertibles.

For $R8\_2\_808$, dimension $186=2\cdot3\cdot31$ and the integral
$pqr$ theorem make any categorification group-theoretical
\cite[Theorem~9.2]{ENOSolvable}. Groups of order $186$ are solvable,
so the category is solvable. The adjoint, type $(1^2,2,6^5)$, and the character involution fixing
all noninvertibles force
$C_2$ descent of type $(1^3,3^{10})$ and dimension ninety-three.
The same prime-dimensional-adjoint argument shows that its universal
grading is trivial. Thus it is $(\Vec_{C_{31}}^\omega)^{C_3}$ with
nontrivial order-three action. Uniqueness of the central $C_3$
preserves it under the outer involution. The combined group is $S_3$,
because its embedded representation subcategory has type $(1^2,2)$.
Every homomorphism $S_3\to\operatorname{Aut}(C_{31})$ kills its normal
$C_3$, contradicting that action.

For the commutative ring $R8\_2\_181$, the neutral category has
dimension twenty-one, type $(1^3,3^2)$ and trivial universal grading.
Solvability makes it $(\Vec_{C_7}^\omega)^{C_3}$, with multiplication
by two or four on $C_7$. The action on
$H^3(C_7,\C^\times)=C_7$ is multiplication by four or two, so its
invariant associator class is zero. Coprime action cohomology then
identifies the neutral category with $\Rep(C_7\rtimes C_3)$.
Its central $C_3$ has unique lifts. The three odd simples have free,
identical left and right $C_3$ actions, so the odd bimodule fixes these
central lifts pointwise. Relative-center equivariantization
\cite{GNN} and $H^2(C_2,\C^\times)=0$ give six central invertible
lifts upstairs. The unique order-three subgroup has trivial
self-braiding, inherited from the neutral $\Rep(C_3)$, and is an
embedded Tannakian subcategory. Its de-equivariantization has the
Tambara--Yamagami rules on $C_7$. The induced $C_3$ symmetry must
preserve its nondegenerate symmetric bicharacter
$\chi(a,b)=\exp(2\pi i kab/7)$, $k\ne0$. Multiplication by two or
four cannot do so: neither square is one modulo seven. This proves
the contradiction without a pivotal or a fermionic-transparency
assumption.
\end{proof}

\subsection{A dimension-twenty-eight adjoint criterion}
\begin{lemma}\label{lem:dim28}
A nonpointed integral fusion category of dimension twenty-eight has
universal grading $C_2$ and adjoint fusion ring $\Rep(D_{14})$.
\end{lemma}
\begin{proof}
Jordan--Larson's theorem \cite[Theorem~1.1]{JordanLarson} applies with
$p=7,q=2$ and $7\nmid q+1$. Write the category as
$\mathcal C(G,\omega,H,\psi)$, where $|G|=28$.
Sylow theory gives $G=C_7\rtimes P$, with $P=C_4$ or $C_2^2$;
the action is trivial or inversion. The possible subgroup orders are
$1,2,4,7,14,28$. For orders one and seven, or a central subgroup of
order two, the cyclic stabilizers give a pointed category.
For a noncentral subgroup $H$ of order two, its normal closure is
$C_7\rtimes H=D_{14}$, of index two; support gives a faithful $C_2$
grading. For $|H|=14$, $H$ is normal and support again gives such a
grading. An abelian $H=C_{14}$ would give a pointed category; otherwise
$H=D_{14}$, whose Schur multiplier is trivial. For $H=G$ one obtains
$\Rep(G)$; a nonabelian group of order twenty-eight has center $C_2$
and quotient $D_{14}$. In each of these nonpointed cases the neutral
fusion ring is $\Rep(D_{14})$.

It remains to justify this ring identification for $|H|=4$, and also
the cocycle reduction used for a reflection complement above.
Conjugate $H$ to $P$. The positive-degree cohomology of $C_7$ is
seven-primary, so the coprime Lyndon--Hochschild--Serre sequence has
no mixed positive-degree terms. Since $\omega|_P=d\psi$, the
associator class is represented by a $P$-invariant class on $C_7$.
Choose a normalized $\mu_7$ representative, average over $P$, and
raise to the inverse of $|P|$ modulo seven. This gives a strictly
invariant representative $\alpha$. On $C_7\rtimes P$ use
\[
\omega_\alpha((n_1,p_1),(n_2,p_2),(n_3,p_3))
 =\alpha(n_1,p_1n_2,p_1p_2n_3).
\]
It is identically one on $P$. The remaining $2$-cocycle $\psi$ is
removed by its inflation along $G\to P$. Thus the category is
$(\Vec_{C_7}^{\alpha})^P$ with strict permutation action.
A trivial action is pointed. Otherwise its kernel $C_2$ acts strictly
trivially and its characters grade the category, with neutral
$(\Vec_{C_7}^{\alpha})^{C_2}$. This has the $\Rep(D_{14})$ fusion
ring; the same argument with a reflection complement supplies the
previous order-two case.

Finally this neutral ring is adjoint to itself. For its three
dimension-two objects, indexed by nonzero elements modulo inversion,
$X_i^2=1+\varepsilon+X_{2i}$, and doubling cycles the three indices.
The adjoint of the full category is contained in the displayed neutral
component and contains its adjoint, so it equals that component.
Consequently the universal grading is exactly $C_2$.
\end{proof}

\begin{proposition}\label{prop:dim56}
There is no fusion category with dimension fifty-six, type
$(1^4,2,4^3)$ and trivial universal grading. In particular
$R8\_2\_694$ is generally excluded.
\end{proposition}
\begin{proof}
Solvability and trivial universal grading force a last equivariantization.
A $C_7$ equivariantization would supply seven invertibles, so the last
group is $C_2$. Its descent $\mathcal D$ has dimension twenty-eight
and is integral by the orbit/stabilizer dimension formula. It cannot
be pointed, since a $C_2$ equivariantization of a pointed category has
all dimensions at most two. Lemma~\ref{lem:dim28} therefore applies.
The outer involution preserves $\mathcal D_{\rm ad}$ and permutes
its three dimension-two simples, fixing at least one. Its cyclic
stabilizer has trivial projective multiplier, so its two characters
produce two distinct dimension-two equivariant simples. The proposed
category has only one. This is a contradiction.
\end{proof}
